\documentclass[11pt]{article}
\usepackage[margin=0.8in]{geometry}
\usepackage[T1]{fontenc}
\usepackage{lmodern}
\usepackage{hyperref}
\hypersetup{colorlinks=true, linkcolor=blue, urlcolor=blue}
\setlength{\parindent}{0pt}
\setlength{\parskip}{0.45em}

\title{General Codex Output}
\author{Codex}
\date{2026-06-22}

\begin{document}
\maketitle

\section*{Available Root Documents}

This Overleaf project is organized as a multi-document output buffer. Compile
one of the root \texttt{.tex} files below depending on which PDF you want.

\begin{itemize}
\item \texttt{d0\_recreated\_analysis\_note.tex}: AN-shaped recreated
\(D^0\) UPC analysis note assembled from reproduced artifacts.
\item \texttt{d0\_reproduction\_workflow.tex}: barebones workflow/timeline for
the reproduced \(D^0\) mass-peak chain from MINIAOD to selected \(D^0\)
candidates and mass plots.
\item \texttt{dplus\_mass\_peak\_search.tex}: live \(D^+ \to K\pi\pi\)
peak-search report, separate from the \(D^0\) AN reproduction.
\end{itemize}

\section*{Project Layout}

Top-level \texttt{.tex} files are intended to be compile roots. Shared figures,
tables, manifests, project notes, and local build outputs live under
\texttt{support/}.

\end{document}
